% Módulo sin preámbulo para manual_usuario_toolbox_chaos.tex.
% Contraste funcional: SYSTEM_REGISTRY, SystemParameterPanel, tab_controls,
% custom_system_tab y sprott_explorer_tab del árbol local de Toolbox Chaos.

\section{Protocolos de simulación y registro experimental}
\label{sec:recetario-ampliado}

Este recetario convierte una secuencia de controles en un experimento que otra
persona puede repetir. Cada receta separa siete componentes: la pregunta, el
modelo y sus valores exactos, la ruta por la interfaz, la salida esperada, la
lectura de la gráfica, una prueba de refinamiento y el límite de la evidencia.
El orden organiza la procedencia del experimento. Cuando se guarda primero una imagen y se reconstruye el
contrato después, es fácil omitir el estado inicial, el transitorio o una
decisión de muestreo que cambió el resultado.

Los nombres de pestañas y campos se reproducen literalmente y se destacan
tipográficamente; las acciones visibles se presentan como botones. Los valores
indicados pertenecen a la implementación actual. Al trabajar con otra versión,
anote cualquier diferencia en la bitácora en vez de ajustar silenciosamente el
protocolo.

\begin{cajaazul}{Regla de trabajo para todo el recetario}
Asigne un identificador a la corrida antes de calcular, cambie una sola familia
de controles durante la prueba de refinamiento y conserve por separado la
salida base y la refinada. En las pestañas generales, \boton{Guardar gráfica...}
guarda la figura. Complete también la plantilla de bitácora con todos los
parámetros de la
sección~\ref{sec:bitacora-reproducible}.
\end{cajaazul}

\subsection{Funciones de las trece pestañas}

La tabla siguiente permite elegir la herramienta por la pregunta que se desea
responder. También señala el límite operativo que debe quedar escrito en el
informe.

\begin{longtable}{@{}>{\raggedright\arraybackslash}p{0.16\linewidth}>{\raggedright\arraybackslash}p{0.22\linewidth}>{\raggedright\arraybackslash}p{0.24\linewidth}>{\raggedright\arraybackslash}p{0.22\linewidth}@{}}
\toprule
Pestaña & Pregunta operativa & Acción o producto principal & Límite que debe registrarse \\
\midrule
\endhead
\campo{Atractor 3D} & ¿Qué geometría muestra una trayectoria de tres componentes? & \boton{Generar atractor 3D} y \boton{Guardar gráfica...}. & El análisis de caos, atracción y estabilidad numérica requiere diagnósticos complementarios. \\
\campo{Retratos 2D} & ¿Qué estructura conservan las proyecciones por pares? & Simulación local o \boton{Usar última trayectoria}; panel de proyecciones. & Una proyección puede superponer estados distintos y ocultar otras coordenadas. \\
\campo{Series temporales} & ¿Cómo cambian las variables con el tiempo o la iteración? & Simulación local o reutilización de la última trayectoria. & La ventana finita puede mezclar transitorio y régimen posterior. \\
\campo{Comparar metodos} & ¿Qué cambia al usar Euler, Heun o RK4 con el mismo paso? & \boton{Comparar integradores}. & En caos, compare escala global, estructura de fase y diagnósticos declarados durante tiempos largos. \\
\campo{Espectro} & ¿Qué bandas o frecuencias aparecen en la muestra? & \campo{PSD de Welch (recomendado)} o \campo{Espectro de amplitud}. & Los máximos dependen del transitorio, la duración y la frecuencia de muestreo. \\
\campo{Lyapunov} & ¿Cómo evolucionan estimaciones finitas de separación local? & Espectro QR--Benettin con RK4 fijo. & Admite flujos ODE enteros registrados de tres dimensiones; una conclusión requiere la secuencia de refinamiento y convergencia del signo. \\
\campo{Bifurcación} & ¿Cómo cambia una observable al barrer un parámetro? & Barrido con transitorio, ventana útil y continuación opcional. & La resolución, el sentido del barrido y la clasificación numérica condicionan la figura. \\
\campo{Cuenca de atracción} & ¿Qué destino calcula cada condición inicial de una malla? & Clasificación en el plano \((x_0,y_0)\) con \(z_0\) fijo. & La certificación de destinos globales y fronteras exactas requiere análisis de cobertura y convergencia. \\
\campo{Autovalores} & ¿Qué predice la linealización cerca de un equilibrio? & Equilibrios y autovalores en el plano complejo. & La dinámica alejada del equilibrio requiere trayectorias, cuencas y análisis globales complementarios. \\
\campo{Coexistencia} & ¿Qué ocurre con varias condiciones iniciales y los mismos parámetros? & Simulación de uno o todos los casos catalogados. & La clasificación de los destinos posibles requiere una cobertura de condiciones iniciales y sus diagnósticos. \\
\campo{Manuales} & ¿Qué fundamento o procedimiento hace falta consultar? & Tres manuales integrados: usuario, teoría y Explorador Sprott. & El registro de la corrida concreta aporta la procedencia experimental. \\
\campo{Crear sistema} & ¿Qué trayectoria produce una definición algebraica propia? & Validación restringida, simulación y archivo JSON. & Equilibrios, Jacobiano, caos, atracción global y cuencas requieren los flujos de análisis y curación correspondientes. \\
\campo{Explorador Sprott} & ¿Qué sistema representa un código compacto y cómo se explora? & Decodificación, simulación, búsqueda, estilo, galería y exportaciones. & Los filtros rápidos producen candidatos; la certificación de caos o atracción requiere diagnósticos adicionales. \\
\bottomrule
\end{longtable}

\subsection{Protocolo 1. Trayectoria de Lorenz en varias pestañas}

\begin{cajaazul}{Trayectoria de Lorenz y diagnósticos}
¿La geometría de dos lóbulos, la irregularidad temporal, el contenido espectral
y la estimación de Lyapunov describen de manera compatible una misma corrida de
Lorenz?
\end{cajaazul}

\paragraph{Parámetros de Lorenz.}
Use \campo{Lorenz [listo]}, \(\sigma=10\), \(\rho=28\),
\(\beta=2.666667\), \(x(0)=y(0)=z(0)=0.1\),
\campo{Runge--Kutta 4 [C]}, \campo{dt}=0.01 y
\campo{Tiempo total T}=40. Para la PSD descarte 5 unidades de tiempo y deje
ambos límites de frecuencia en cero. Para Lyapunov use
\campo{dt integrador}=0.01, \campo{Burn-in time}=5,
\campo{Tiempo final}=40 y \campo{QR cada N pasos}=10.

\paragraph{Controles de la simulación.}
\begin{enumerate}
  \item Abra \campo{Atractor 3D}, capture el contrato anterior y pulse
  \boton{Generar atractor 3D}. Guarde la figura como
  \archivo{r01_lorenz_atractor_base.png}.
  \item Abra \campo{Retratos 2D}, seleccione también
  \campo{Lorenz [listo]} y pulse \boton{Usar última trayectoria}. Conserve los
  planos \(x\)-\(y\), \(x\)-\(z\) y \(y\)-\(z\); guarde
  \archivo{r01_lorenz_retratos_base.png}.
  \item En \campo{Series temporales}, seleccione Lorenz y pulse
  \boton{Usar última trayectoria}. Guarde
  \archivo{r01_lorenz_series_base.png}.
  \item En \campo{Espectro}, seleccione
  \campo{PSD de Welch (recomendado)}, escriba 5 en
  \campo{Descartar transitorio (s)}, deje los límites automáticos en cero y
  pulse \boton{Usar última trayectoria}. Guarde
  \archivo{r01_lorenz_psd_base.png}.
  \item En \campo{Lyapunov}, seleccione Lorenz, capture los cuatro valores
  indicados, pulse \boton{Calcular exponentes de Lyapunov} y guarde
  \archivo{r01_lorenz_lyapunov_base.png}.
\end{enumerate}

\begin{cajaverde}{Geometría, series y Lyapunov de Lorenz}
La vista 3D debe organizarse alrededor de dos lóbulos; los retratos 2D muestran
proyecciones diferentes de los mismos estados y las series alternan episodios
asociados a ambos lóbulos. Una PSD distribuida en varias frecuencias es
compatible con una señal de varias componentes o con modulación irregular. En Lyapunov importa la curva
de convergencia además del valor final: registre los tres exponentes mostrados
y si el máximo conserva su signo durante la parte tardía de la ventana.
\end{cajaverde}

\figuradidactica{systems/lorenz_phase_timeseries.png}{Lectura conjunta de Lorenz: la geometría tridimensional identifica los dos lóbulos y las series temporales muestran cuándo la trayectoria cambia de uno a otro. Los ejes y el pie contienen la información descriptiva y mantienen una composición gráfica limpia.}

\paragraph{Refinamiento del paso temporal.}
Repita la trayectoria con \campo{dt}=0.005 y \campo{Tiempo total T}=40;
regenere las tres representaciones desde esa nueva corrida. Para Lyapunov use
\campo{dt integrador}=0.005, \campo{Burn-in time}=5,
\campo{Tiempo final}=80 y \campo{QR cada N pasos}=10. Compare extensión,
rangos, lóbulos, bandas espectrales y tendencia de los exponentes. Use esos
indicadores para comparar las trayectorias tardías.

\begin{cajaalerta}{Diagnósticos para caos y atracción en Lorenz}
La identificación de caos y atracción requiere varios diagnósticos. La forma de mariposa y una PSD ancha orientan el análisis, mientras que un exponente máximo positivo aporta evidencia numérica dentro del
contrato usado. Una demostración matemática y la convergencia entre plataformas
requieren sus desarrollos y verificaciones correspondientes.
\end{cajaalerta}

\paragraph{Archivos y metadatos de Lorenz.}
Conserve las cinco figuras base, sus versiones con sufijo
\archivo{_dt0005}, y \archivo{r01_lorenz_bitacora.txt}. En la bitácora anote
versión, plataforma, parámetros, estado inicial, método, ambos pasos, tiempos,
transitorio de la PSD, controles de Lyapunov, nombres de salida y los valores
numéricos reportados por la interfaz.

\subsection{Protocolo 2. Mapa logístico y diagrama de bifurcación}

\begin{cajaazul}{Bifurcaciones del mapa logístico}
¿Qué cambios cualitativos aparecen en los valores tardíos de
\(x_{n+1}=r x_n(1-x_n)\) cuando \(r\) recorre el intervalo de 2.5 a 4.0?
\end{cajaazul}

\paragraph{Parámetros del mapa logístico.}
En \campo{Bifurcación} seleccione \campo{Logistico [listo]},
\(x(0)=0.2\), \campo{Parámetro de bifurcación}=\campo{r},
\campo{Variable observada}=\campo{x}, \campo{Mínimo}=2.5,
\campo{Máximo}=4.0, \campo{N (Parámetros)}=350,
\campo{dt integrador}=1, \campo{Transitorio descartado}=600,
\campo{Tiempo útil}=500 y \campo{Máx cruces por param}=250.
Mantenga desactivado \campo{Usar continuación} en la corrida base.

\paragraph{Controles del diagrama de bifurcación.}
\begin{enumerate}
  \item Confirme que el selector del mapa deshabilita el integrador: una
  iteración pertenece al modelo; el paso temporal corresponde a una EDO.
  \item Capture los valores del barrido y pulse
  \boton{Calcular bifurcación}.
  \item Lea primero el eje horizontal y confirme que cubre 2.5--4.0; después
  examine cuántos valores tardíos aparecen para cada parámetro.
  \item Guarde \archivo{r02_logistico_bif_base.png}.
\end{enumerate}

\begin{cajaverde}{Ramas y bandas del mapa logístico}
La figura debe pasar de una rama a duplicaciones sucesivas, regiones densas y
ventanas más ordenadas. Una columna representa una muestra tardía para un valor
de \(r\); su densidad vertical registra la multiplicidad de valores y puede estar limitada por
\campo{Máx cruces por param}. Distinga una ventana visible de la afirmación de
que se conoce su periodo exacto.
\end{cajaverde}

\figuradidactica{logistic_bifurcation.png}{Barrido reproducible del mapa logístico. Una sola altura por parámetro indica un estado asintótico fijo; dos, cuatro y más alturas revelan duplicaciones de periodo; las bandas densas deben contrastarse con series y refinamiento.}

\paragraph{Resolución del barrido paramétrico.}
Repita sin continuación con \campo{N (Parámetros)}=700,
\campo{Transitorio descartado}=1000, \campo{Tiempo útil}=800 y
\campo{Máx cruces por param}=500. Guarde
\archivo{r02_logistico_bif_refinada.png}. Después, si desea estudiar dependencia
de rama, active \campo{Usar continuación} y guarde esa corrida con el sufijo
\archivo{_continuacion}; manténgala como una corrida separada de la prueba de resolución.

\begin{cajaalerta}{Resolución y parámetros críticos del mapa logístico}
La caracterización de cada banda y la localización de parámetros críticos requieren
resolver el barrido con precisión suficiente. La resolución horizontal, el
transitorio y la cantidad de valores guardados condicionan las estructuras
pequeñas.
\end{cajaalerta}

\paragraph{Archivos y metadatos del mapa logístico.}
Guarde las dos figuras principales, la comparación opcional de continuación y
\archivo{r02_logistico_bitacora.txt}. Registre ecuación, \(x(0)\), intervalo,
número de parámetros, transitorio, ventana útil, límite por columna, estado de
continuación y dirección del barrido.

\subsection{Protocolo 3. Mapa de Hénon y retrato de fase}

\begin{cajaazul}{Retrato de fase del mapa de Hénon}
¿Cómo se organiza en el plano \(x\)-\(y\) la órbita del mapa de Hénon y qué
aspectos permanecen al aumentar el número de iteraciones?
\end{cajaazul}

\paragraph{Parámetros del mapa de Hénon.}
Use \campo{Henon [listo]}, \(a=1.4\), \(b=0.3\),
\(x(0)=0.1\), \(y(0)=0.1\), \campo{dt}=1 y
\campo{Tiempo total T}=5000. El campo de método debe quedar deshabilitado por
tratarse de un mapa. El estado 2D queda definido por las dos primeras condiciones iniciales.

\paragraph{Controles de iteración y retrato.}
\begin{enumerate}
  \item Abra \campo{Retratos 2D}, seleccione Hénon y confirme los valores.
  \item En \campo{Proyecciones 2D}, conserve \campo{Plano x-y}.
  \item Pulse \boton{Generar retratos 2D} y guarde
  \archivo{r03_henon_xy_base.png}.
  \item Abra \campo{Series temporales}, seleccione Hénon y pulse
  \boton{Usar última trayectoria}; guarde
  \archivo{r03_henon_series_base.png}.
\end{enumerate}

\begin{cajaverde}{Geometría plegada e iteraciones de Hénon}
El plano debe mostrar un conjunto plegado y delgado, con estructura distinta de una curva
cerrada simple. Las series de \(x_n\) y \(y_n\) permiten comprobar que cada
punto procede de la iteración anterior. La pestaña 3D trabaja con tres
coordenadas, mientras este sistema registra dimensión dos; use el plano \(x\)-\(y\).
\end{cajaverde}

\figuradidactica{systems/henon_phase_timeseries.png}{El retrato de Hénon conserva la geometría de las iteraciones después del transitorio, mientras las series separan el orden temporal que se pierde en la nube del plano.}

\paragraph{Longitud de la iteración.}
En un mapa, refine la longitud de la iteración. Repita con
\campo{Tiempo total T}=10000 y los demás valores intactos. Compare el soporte,
los rangos y las zonas de mayor densidad visual. Como prueba separada de
sensibilidad, use \(x(0)=0.100001\), \(y(0)=0.1\); etiquete esa salida como
perturbación inicial y regístrela como una prueba de sensibilidad distinta del refinamiento temporal.

\begin{cajaalerta}{Diagnósticos de invariancia y caos en Hénon}
La forma plegada aporta geometría proyectada. La dimensión fractal, la ergodicidad,
la invariancia y el caos requieren sus diagnósticos correspondientes. El panel incluye el arranque de la órbita; registre el intervalo de análisis seleccionado para la interpretación posterior.
\end{cajaalerta}

\paragraph{Archivos y metadatos del mapa de Hénon.}
Conserve \archivo{r03_henon_xy_base.png},
\archivo{r03_henon_xy_n10000.png}, las series y, si se ejecutó, la perturbación.
Registre que \(T\) se interpreta como número de iteraciones porque \(dt=1\),
además de parámetros, estado inicial y cantidad total de puntos.

\subsection{Protocolo 4. Sistema de Rössler e integradores}

\begin{cajaazul}{Sensibilidad numérica del sistema de Rössler}
¿Qué rasgos globales del flujo de Rössler son estables frente a Euler, Heun y
RK4, y en qué momento dejan de coincidir sus trayectorias?
\end{cajaazul}

\paragraph{Parámetros del sistema de Rössler.}
En \campo{Comparar metodos} seleccione \campo{Rossler [listo]},
\(a=0.2\), \(b=0.2\), \(c=5.7\),
\((x(0),y(0),z(0))=(0.1,0,0)\), \campo{dt}=0.01 y
\campo{Tiempo total T}=80. Active \campo{Euler explícito},
\campo{Heun / Euler mejorado (RK2)} y \campo{Runge--Kutta 4}; asigne un color
distinto a cada método.

\paragraph{Controles de integración.}
\begin{enumerate}
  \item Confirme que los tres integradores están activos y pulse
  \boton{Comparar integradores}.
  \item Guarde \archivo{r04_rossler_metodos_dt001.png}.
  \item Abra \campo{Retratos 2D}, capture el mismo modelo con
  \campo{Runge--Kutta 4 [C]}, \(dt=0.01\), \(T=80\), genere los retratos y
  guarde \archivo{r04_rossler_retratos_rk4.png}.
\end{enumerate}

\begin{cajaverde}{Trayectorias e integradores de Rössler}
Las series deben mostrar oscilaciones que crecen alrededor de una región y una
excursión de retorno; en el plano aparece la geometría espiral característica.
Las curvas pueden separarse en fase con el tiempo. Compare amplitudes, rangos,
periodo aproximado de las vueltas y presencia de la excursión, junto con
la convergencia de las propiedades reportadas.
\end{cajaverde}

\figuradidactica{systems/rossler_phase_timeseries.png}{Rössler en espacio de fases y en el tiempo. La espiral casi plana explica por qué distintas proyecciones pueden parecer regulares aunque las excursiones de la tercera coordenada introduzcan otra escala dinámica.}

\paragraph{Convergencia al reducir el paso.}
Repita la comparación con \campo{dt}=0.005 y \campo{Tiempo total T}=80.
Guarde \archivo{r04_rossler_metodos_dt0005.png}. Compruebe si Heun y RK4
conservan durante más tiempo la escala y la fase, y registre si Euler presenta
sesgo o inestabilidad. Conserve cada salida inesperada como parte del
diagnóstico.

\begin{cajaalerta}{Convergencia de los integradores de Rössler}
La estimación del error global requiere comparar cuantitativamente los métodos y sus
refinamientos. Una divergencia tardía orienta la revisión del horizonte y del paso. Hace falta
una secuencia de pasos y observables cuantitativas para sostener convergencia.
\end{cajaalerta}

\paragraph{Archivos y metadatos de Rössler.}
Guarde las dos comparaciones, el retrato RK4 y
\archivo{r04_rossler_bitacora.txt}. Anote versión, ecuaciones de referencia,
parámetros, estado inicial, métodos activados, colores, pasos, horizonte y el
primer intervalo temporal en que las curvas dejan de ser visualmente próximas.

\subsection{Protocolo 5. Circuito de Chua y condición inicial}

\begin{cajaazul}{Condiciones iniciales y simetría de Chua}
¿Dos condiciones iniciales simétricas, con exactamente los mismos parámetros de
Chua, producen trayectorias reflejadas y una geometría de doble scroll durante
el horizonte simulado?
\end{cajaazul}

\paragraph{Parámetros del circuito de Chua.}
En \campo{Atractor 3D} elija \campo{Chua}. Use
\(\alpha=15.6\), \(\beta=28\), \(m_0=-1.143\), \(m_1=-0.714\), con
\((0.1,0,0)\) en la primera corrida y \((-0.1,0,0)\) en la segunda.
Seleccione RK4, \campo{dt}=0.01 y \campo{Tiempo T}=80.

\paragraph{Controles de la simulación de Chua.}
\begin{enumerate}
  \item Introduzca los parámetros y la condición inicial \((0.1,0,0)\), pulse
  \boton{Generar atractor 3D} y guarde \archivo{r05_chua_ic_positiva.png}.
  \item Cambie únicamente la condición inicial a \((-0.1,0,0)\), repita la
  simulación y guarde \archivo{r05_chua_ic_negativa.png}.
  \item Abra \campo{Series temporales} y conserve las componentes de ambas
  corridas con nombres que identifiquen claramente el signo del inicio.
  \item Compare \((x,y,z)\) de una corrida con \((-x,-y,-z)\) de la otra. La
  simetría del campo vectorial permite comprobar numéricamente el reflejo sin
  confundirlo con dos sistemas diferentes.
\end{enumerate}

\begin{cajaverde}{Simetría numérica en las trayectorias de Chua}
Las dos corridas deben ser finitas, recorrer una geometría de doble scroll y
aparecer como reflejos dentro del error numérico. Lea la comparación como una
verificación de simetría y de sensibilidad de la trayectoria concreta al signo
del inicio. La existencia de dos destinos asintóticos distintos requiere una
exploración de cuencas y condiciones iniciales.
\end{cajaverde}

\figuradidactica{systems/chua_phase_timeseries.png}{Chua con trayectoria tridimensional y componentes temporales. La figura permite comprobar doble scroll, finitud y escalas; la pareja de condiciones iniciales de la receta debe analizarse además mediante el error de simetría.}

\paragraph{Sensibilidad a la condición inicial.}
Repita ambas corridas con \campo{dt}=0.005 y \campo{Tiempo T}=80. Después
ejecute una tercera pareja con \campo{dt}=0.01 y \campo{Tiempo T}=160 para
separar el efecto del paso del efecto del horizonte. Compare la región visitada,
la simetría y las estadísticas globales como criterios de comparación a tiempos
tardíos.

\begin{cajaalerta}{Cuencas y coexistencia en el circuito de Chua}
Esta reproducción muestra dos trayectorias simétricas para el protocolo
declarado. La caracterización de caos, coexistencia de destinos asintóticos y fronteras de cuencas requiere ampliar las condiciones iniciales y los diagnósticos.
\end{cajaalerta}

\paragraph{Archivos y metadatos de Chua.}
Conserve las figuras y series de las dos corridas, las pruebas de refinamiento y
\archivo{r05_chua_bitacora.txt}. Incluya parámetros, ambas condiciones
iniciales, método, \(dt\), \(T\), norma del error de simetría y nombres de
archivo.

\subsection{Protocolo 6. Duffing--Ueda y muestreo del forzamiento}

\begin{cajaazul}{Muestreo del forzamiento Duffing--Ueda}
¿El paso temporal resuelve suficientemente el periodo del forzamiento y se
conservan la envolvente y las bandas espectrales de la respuesta Duffing--Ueda?
\end{cajaazul}

\paragraph{Parámetros de Duffing--Ueda.}
En \campo{Series temporales} seleccione \campo{Duffing-Ueda [listo]},
\(\delta=0.2\), \(\alpha=-1\), \(\beta=1\), \(\gamma=0.3\),
\(\omega=1.2\), \((x(0),y(0),\theta(0))=(0.1,0,0)\),
\campo{Runge--Kutta 4 [C]}, \campo{dt}=0.01 y
\campo{Tiempo total T}=200. El periodo de forzamiento es
\[
  P_f=\frac{2\pi}{\omega}=5.235988\ldots,
\]
por lo que la malla base contiene aproximadamente 523.6 muestras por periodo.

\paragraph{Controles de integración y muestreo.}
\begin{enumerate}
  \item Pulse \boton{Generar series temporales} y guarde
  \archivo{r06_duffing_series_dt001.png}.
  \item Abra \campo{Retratos 2D}, seleccione Duffing--Ueda y pulse
  \boton{Usar última trayectoria}. Conserve el plano \(x\)-\(y\) y guarde
  \archivo{r06_duffing_xy_dt001.png}.
  \item Abra \campo{Espectro}, seleccione
  \campo{PSD de Welch (recomendado)}, escriba 50 en
  \campo{Descartar transitorio (s)}, deje los límites de frecuencia en cero y
  pulse \boton{Usar última trayectoria}. Guarde
  \archivo{r06_duffing_psd_dt001.png}.
\end{enumerate}

\begin{cajaverde}{Respuesta temporal y espectral de Duffing--Ueda}
La tercera variable debe crecer como fase extendida porque
\(\dot\theta=\omega\); la fuerza usa \(\cos\theta\), de modo que el fenómeno
físico asociado sigue siendo periódico. Lea en \(x(t)\) y \(y(t)\) la amplitud,
los cambios de pozo y la regularidad de los ciclos. En la PSD, la frecuencia de
forzamiento en ciclos por unidad de tiempo es
\(f_f=\omega/(2\pi)\approx0.190986\); anote bandas cercanas y sus armónicos sin
suponer que todo máximo procede directamente del forzamiento.
\end{cajaverde}

\figuradidactica{systems/duffing_ueda_phase_timeseries.png}{Respuesta forzada de Duffing--Ueda. El retrato proyectado resume la región visitada y las series permiten reconocer el periodo del forzamiento, el transitorio y los cambios al refinar el paso.}

\paragraph{Resolución temporal del forzamiento.}
Repita con \campo{dt}=0.005 y \campo{Tiempo total T}=200; ahora hay unas
1047.2 muestras por periodo. Mantenga el mismo descarte de 50 y compare rangos,
envolvente del plano \(x\)-\(y\) y posiciones de los máximos espectrales. La
interfaz general analiza las series temporales directamente. Una sección de
Poincaré requiere un muestreo estroboscópico dedicado, una vez por periodo.

\begin{cajaalerta}{Caos y convergencia en Duffing--Ueda}
Resolver el forzamiento con muchas muestras permite estudiar la respuesta temporal; el caos y la convergencia global requieren diagnósticos adicionales. Una rampa de \(\theta\) corresponde al avance de fase del forzamiento y
la fase extendida que usa la formulación autónoma correspondiente.
\end{cajaalerta}

\paragraph{Archivos y metadatos de Duffing--Ueda.}
Guarde las seis figuras base/refinadas y
\archivo{r06_duffing_muestreo.txt}. Registre \(P_f\), \(f_f\), muestras por
periodo, parámetros, estado inicial, método, paso, horizonte, transitorio de la
PSD y si se observaron cambios de pozo.

\subsection{Protocolo 7. Mackey--Glass e historia retardada}

\begin{cajaazul}{Historia retardada de Mackey--Glass}
¿La oscilación de Mackey--Glass y su contenido espectral permanecen al refinar
la interpolación temporal de una historia con retardo \(\tau=17\)?
\end{cajaazul}

\paragraph{Parámetros de Mackey--Glass.}
En \campo{Series temporales} seleccione \campo{Mackey-Glass [listo]},
\(\beta=0.2\), \(\gamma=0.1\), \(n=10\), \(\tau=17\),
\(x(0)=1.2\), \campo{Runge--Kutta 4 [C]}, \campo{dt}=0.1 y
\campo{Tiempo total T}=500. Deje en cero los dos campos iniciales mostrados con
guion: el backend construye la historia previa como constante igual a
\(x(0)\). Para RK4 usa interpolación cúbica causal por tramos.

\paragraph{Controles de historia y ventana.}
\begin{enumerate}
  \item Pulse \boton{Generar series temporales} y guarde
  \archivo{r07_mackey_series_dt01.png}.
  \item Identifique los tres canales devueltos por la implementación: valor
  actual, valor retardado y derivada. Los guiones de la interfaz identifican
  los campos de una historia con retardo.
  \item Abra \campo{Espectro}, seleccione
  \campo{PSD de Welch (recomendado)}, establezca
  \campo{Descartar transitorio (s)}=100, deje los límites automáticos y pulse
  \boton{Usar última trayectoria}. Guarde
  \archivo{r07_mackey_psd_dt01.png}.
\end{enumerate}

\begin{cajaverde}{Canales temporales y espectro de Mackey--Glass}
Después de la historia inicial deben aparecer oscilaciones acotadas cuya
amplitud y separación entre máximos pueden variar. Compare el canal actual con
el retardado para reconocer el desplazamiento de \(\tau\). La PSD resume la
ventana posterior a \(t=100\); el análisis de la ecuación diferencial funcional complementa las bandas.
\end{cajaverde}

\figuradidactica{systems/mackey_glass_phase_timeseries.png}{Mackey--Glass requiere leer una ventana larga: la proyección organiza estados retardados y la serie temporal muestra si el horizonte contiene suficientes oscilaciones para sostener la interpretación.}

\paragraph{Longitud de la ventana de observación.}
Repita con \campo{dt}=0.05 y \campo{Tiempo total T}=500, sin cambiar historia,
parámetros ni método. Vuelva a descartar 100 en la PSD. Compare rangos, tiempos
aproximados de máximos y ubicación de bandas. Si difieren materialmente, pruebe
\(dt=0.025\) antes de interpretar la corrida.

\begin{cajaalerta}{Diagnósticos de historias y atracción en Mackey--Glass}
Una señal irregular requiere diagnósticos de caos en el espacio de historias, y la atracción global exige explorar varias historias iniciales. La representación de tres canales es
una vista de una dinámica con memoria; la EDO autónoma tridimensional ordinaria
requiere una formulación de estado distinta.
\end{cajaalerta}

\paragraph{Archivos y metadatos de Mackey--Glass.}
Conserve las series y PSD para ambos pasos y
\archivo{r07_mackey_bitacora.txt}. Incluya la historia constante, \(x(0)\),
\(\beta,\gamma,n,\tau\), interpolación asociada a RK4, pasos, horizonte,
transitorio, versión y cualquier advertencia de memoria o no finitud.

\subsection{Protocolo 8. Lorenz--96 y proyecciones de ocho variables}

\begin{cajaazul}{Proyecciones de ocho variables de Lorenz--96}
¿Qué evolución muestran las tres primeras componentes de Lorenz--96 cuando una
perturbación pequeña rompe el estado uniforme de un anillo con \(J=8\)?
\end{cajaazul}

\paragraph{Parámetros de Lorenz--96.}
Use \campo{Lorenz-96 [listo]}, \(F=8\), \(J=8\),
\((X_1(0),X_2(0),X_3(0))=(8.01,8,8)\),
\campo{Runge--Kutta 4 [C]}, \campo{dt}=0.01 y
\campo{Tiempo total T}=40. El backend inicializa las restantes componentes del
anillo con \(F\) y devuelve a la interfaz sólo \(X_1,X_2,X_3\).

\paragraph{Controles de la proyección.}
\begin{enumerate}
  \item Abra \campo{Atractor 3D}, capture los valores y pulse
  \boton{Generar atractor 3D}. Guarde
  \archivo{r08_lorenz96_proyeccion3d_base.png}.
  \item Abra \campo{Retratos 2D}, seleccione Lorenz--96 y pulse
  \boton{Usar última trayectoria}. Guarde
  \archivo{r08_lorenz96_retratos_base.png}.
  \item En \campo{Series temporales}, seleccione Lorenz--96, pulse
  \boton{Usar última trayectoria} y guarde
  \archivo{r08_lorenz96_series_base.png}.
\end{enumerate}

\begin{cajaverde}{Propagación y proyecciones de Lorenz--96}
La perturbación inicial debe propagarse y las tres componentes mostradas dejan
de permanecer exactamente en el nivel uniforme. Lea la vista 3D y los planos
como proyecciones de una trayectoria de ocho variables. La cercanía o cruce de
dos puntos en \((X_1,X_2,X_3)\) aportan una coincidencia proyectada; la igualdad
del estado completo requiere comparar las ocho componentes.
\end{cajaverde}

\figuradidactica{systems/lorenz96_phase_timeseries.png}{Proyección de Lorenz--96. La salida gráfica actual muestra tres componentes; el pie recuerda que el estado integrado contiene ocho variables y presenta una vista parcial del sistema completo.}

\paragraph{Selección de componentes observadas.}
Repita con \campo{dt}=0.005 y \campo{Tiempo total T}=40. Compare rangos,
estructura proyectada y escalas temporales. Como experimento distinto del
refinamiento, puede cambiar \(J\); tal cambio modifica la dimensión del modelo
y requiere una nueva bitácora.

\begin{cajaalerta}{Transporte y extensividad en Lorenz--96}
La proyección de tres componentes describe una vista parcial del estado de
Lorenz--96. El transporte en el anillo y la extensividad del caos requieren cálculos específicos. La etiqueta ``atractor 3D'' se reserva para la
ocho.
\end{cajaalerta}

\paragraph{Archivos y metadatos de Lorenz--96.}
Guarde las seis figuras base/refinadas y
\archivo{r08_lorenz96_bitacora.txt}. Registre \(F\), \(J\), las tres entradas
visibles, la regla de inicialización de las cinco restantes, método, paso,
horizonte y las coordenadas proyectadas.

\subsection{Protocolo 9. Genesio--Tesi en Crear sistema}

\begin{cajaazul}{Definición de Genesio--Tesi en el editor}
¿Puede el editor restringido validar, simular y conservar una definición de
Genesio--Tesi sin ejecutar código Python arbitrario?
\end{cajaazul}

\paragraph{Ecuaciones y parámetros de Genesio--Tesi.}
En \campo{Crear sistema} capture los campos siguientes.

\begin{tabularx}{\linewidth}{>{\raggedright\arraybackslash}p{0.34\linewidth}X}
\toprule
Campo & Valor \\
\midrule
\campo{Nombre} & \archivo{Genesio-Tesi} \\
\campo{Tipo} & \campo{Flujo continuo} \\
\campo{Variables (separadas por coma)} & \archivo{x, y, z} \\
\campo{Parametros (nombre=valor)} & \archivo{a=1.2}, \archivo{b=2.92}, \archivo{c=6.0}, una asignación por línea \\
\campo{Ecuaciones (una por variable)} & \archivo{y}; \archivo{z}; \archivo{-c*x-b*y-a*z+x**2}, una expresión por línea \\
\campo{Estado inicial} & \archivo{0.2, 0.2, 0.2} \\
\campo{Metodo del flujo} & \archivo{rk4} \\
\campo{Paso h} & \archivo{0.005000} \\
\campo{Duracion} & \archivo{500.0000} \\
\campo{Muestras transitorias} & \archivo{5000} \\
\campo{Umbral de divergencia} & \archivo{1000000.00} \\
\bottomrule
\end{tabularx}

\paragraph{Controles de Crear sistema.}
\begin{enumerate}
  \item Pulse \boton{Validar definicion}. Exija el mensaje de definición válida
  antes de simular.
  \item Pulse \boton{Simular y graficar}. Registre estado, pasos completados,
  pasos solicitados y motor mostrado.
  \item Pulse \boton{Guardar JSON} y use el nombre
  \archivo{r09_genesio_tesi_definicion.json}.
  \item Abra \campo{Espectro}, fije
  \campo{PSD de Welch (recomendado)},
  \campo{Descartar transitorio (s)}=25 y pulse
  \boton{Usar última trayectoria}. Guarde
  \archivo{r09_genesio_tesi_psd_base.png}.
\end{enumerate}

\begin{cajaverde}{Trayectoria y PSD de Genesio--Tesi}
La validación debe reconocer tres variables y un flujo. La simulación debe
completar una trayectoria finita o informar explícitamente el punto de fallo.
El gráfico integrado muestra la proyección de las tres variables después de
retirar las muestras transitorias elegidas. El JSON conserva la definición
algebraica; la PSD permite revisar escalas de la trayectoria compartida.
\end{cajaverde}

\figuradidactica{variational_volume_manifolds.png}{Apoyo para interpretar un sistema personalizado: el campo determina la órbita, su variacional describe la deformación local y las direcciones estables e inestables organizan el movimiento próximo. La interpretación integra la simulación con esa lectura variacional.}

\paragraph{Verificación de la definición ingresada.}
Repita con \campo{Paso h}=0.002500, \campo{Duracion}=500 y
\campo{Muestras transitorias}=10000. El cambio de 5000 a 10000 mantiene en 25
unidades de tiempo el tramo retirado. Compare región visitada, rangos y PSD;
guarde la figura refinada con el sufijo \archivo{_h00025}.

\begin{cajaalerta}{Evidencia necesaria y alcance de la reutilización}
Una definición válida satisface el contrato sintáctico y numérico del
motor. La evaluación de caos y atracción global requiere diagnósticos posteriores. En la interfaz actual,
\campo{Retratos 2D} y \campo{Series temporales} comparan la clave de la última
trayectoria con el catálogo y pueden rechazar la clave personalizada. La ruta
verificada para reutilizarla es \campo{Espectro}. Registre la compatibilidad
mostrada por la interfaz antes de declarar la reutilización en los otros botones.
\end{cajaalerta}

\paragraph{Archivos y metadatos de Genesio--Tesi.}
Conserve el JSON, ambas PSD y \archivo{r09_genesio_tesi_bitacora.txt}. El JSON conserva la definición algebraica. Añada a la bitácora el paso, la duración,
las muestras transitorias, el umbral, la versión del motor, el estado y el conteo de pasos.

\subsection{Protocolo 10. Código Sprott y exportación}

\begin{cajaazul}{Código compacto y exportación de Sprott}
¿Puede reconstruirse un mapa 2D desde un código compacto y conservarse la
trayectoria separada de las decisiones de estilo?
\end{cajaazul}

\paragraph{Código y parámetros Sprott.}
Use el código \archivo{EWMWAMMMPMMMM}. En \campo{Exploracion} establezca
\campo{Tipo}=\archivo{map}, \campo{Dimension}=2, \campo{Orden}=2,
\campo{\mbox{Iteraciones}/\allowbreak\mbox{pasos}}=12000, \campo{Transitorio}=1500,
\campo{Tolerancia divergencia}=1000000, \campo{Semilla aleatoria}=1 y
\campo{Codigo}=\archivo{EWMWAMMMPMMMM}. Para un código fijo, semilla,
\campo{Paso h} y \campo{Metodo flujo} documentan el estado visible de esta interfaz; la iteración del mapa queda definida por el código, la semilla y los parámetros. Manténgalos
en 1, 0.01 y \archivo{rk4}.

Use además \campo{Proyeccion}=\archivo{x-y},
\campo{Color por}=\archivo{tiempo}, \campo{Paleta}=\campo{Viridis},
\campo{Fondo}=\campo{blanco}, \campo{Modo dibujo}=\campo{puntos},
\campo{Tamano punto}=0.70, \campo{Opacidad}=0.75,
\campo{Max puntos}=12000, \campo{Bandas}=0 y
\campo{DPI export}=220. Active \campo{Mostrar ejes},
\campo{Mostrar grid} y \campo{Aspect ratio igual}.

\paragraph{Controles de estilo y exportación.}
\begin{enumerate}
  \item En \campo{Codigos}, pegue el código y pulse \boton{Decodificar}.
  Verifique que las ecuaciones reconstruidas sean
  \[
    x_{n+1}=1+y_n-1.2x_n^2,
    \qquad y_{n+1}=0.3x_n.
  \]
  \item Pase a \campo{Exploracion}, capture el contrato anterior y pulse
  \boton{Simular y graficar}. Lea estado, razón y rangos antes de cambiar el
  estilo.
  \item Capture los controles visuales y pulse \boton{Aplicar estilo}. Cambiar
  estilo debe conservar la órbita almacenada y modificar únicamente su representación.
  \item Pulse \boton{Exportar imagen actual},
  \boton{Exportar metadatos JSON} y \boton{Exportar datos CSV}; use el prefijo
  \archivo{r10_sprott_EWMWAMMMPMMMM_base}.
  \item Pulse \boton{Copiar codigo} y \boton{Copiar cita}; pegue ambos textos
  en \archivo{r10_sprott_cita.txt}.
\end{enumerate}

\begin{cajaverde}{Órbita y filtro de Sprott}
El decodificador debe identificar una familia polinomial de mapa 2D y doce
coeficientes después de la letra de familia. La órbita postransitoria debe
formar un conjunto plegado bidimensional; el color por tiempo ayuda a seguir la
secuencia sin modificar las coordenadas. Registre el estado real del filtro. Si
aparece \archivo{candidate_chaotic}, léalo como trayectoria finita, acotada y
no colapsada bajo filtros rápidos.
\end{cajaverde}

\figuradidactica{sprott_map4d_turbo_white.png}{Proyección de una órbita 4D del Explorador Sprott: dos coordenadas ocupan los ejes y una coordenada adicional controla el color. El código, el transitorio y el rango cromático deben guardarse con la imagen.}

\paragraph{Verificación del diagnóstico y la figura.}
Mantenga código y estilo, cambie sólo
\campo{\mbox{Iteraciones}/\allowbreak\mbox{pasos}}=24000 y \campo{Transitorio}=3000, y vuelva a
simular. Exporte con el sufijo \archivo{_n24000}. Compare soporte, rangos,
zonas densas y estado del filtro. Una segunda prueba puramente visual puede
cambiar paleta o color, pero debe reutilizar los mismos datos y llevar un sufijo
\archivo{_estilo}; ese sufijo identifica una variación de estilo sobre los mismos datos.

\begin{cajaalerta}{Diagnósticos de caos en candidatos Sprott}
La etiqueta \archivo{candidate_chaotic}, una imagen compleja y una estimación
rápida de Lyapunov forman una primera capa de evidencia. El caos, la existencia de un atractor, la dimensión
fractal y la robustez paramétrica requieren pruebas adicionales. El Explorador Sprott organiza
candidatos y exporta evidencia para análisis posterior.
\end{cajaalerta}

\paragraph{Archivos y metadatos del código Sprott.}
El paquete mínimo contiene PNG, CSV, JSON, cita, bitácora y sus versiones
refinadas. Compruebe que el código del JSON coincide con el copiado; registre
familia, dimensión, orden, ecuaciones, iteraciones, transitorio, umbral, estado,
razón, rangos, proyección, color, paleta, fondo, modo, punto, opacidad, máximo de
puntos, bandas, DPI, versión y procedencia sintética.

\subsection{Campos de la bitácora experimental}
\label{sec:bitacora-reproducible}

Complete una copia por corrida, incluida cada prueba de refinamiento. Cuando un
campo no aplique, escriba ``no aplica''; registre una nota de procedencia en
cada espacio que distinga una omisión de una decisión consciente.

\begin{longtable}{@{}>{\raggedright\arraybackslash}p{0.28\linewidth}>{\raggedright\arraybackslash}p{0.62\linewidth}@{}}
\toprule
Campo & Registro de la corrida \\
\midrule
\endhead
Identificador único & \rule{0pt}{2.8ex} \\
Receta y pregunta & \rule{0pt}{3.5ex} \\
Fecha y hora local & \rule{0pt}{2.8ex} \\
Toolbox Chaos, versión y commit si se conoce & \rule{0pt}{3.5ex} \\
Sistema operativo y equipo & \rule{0pt}{3.5ex} \\
Pestaña y acción ejecutada & \rule{0pt}{3.5ex} \\
Sistema: etiqueta visible y clave si se conoce & \rule{0pt}{3.5ex} \\
Tipo y dimensión del modelo & \rule{0pt}{2.8ex} \\
Fuente, DOI o procedencia & \rule{0pt}{3.5ex} \\
Parámetros, con orden y cifras & \rule{0pt}{4.5ex} \\
Estado o historia inicial & \rule{0pt}{4.5ex} \\
Método y backend mostrados & \rule{0pt}{3.5ex} \\
Paso \(dt\) o \(h\) y horizonte \(T\) & \rule{0pt}{3.5ex} \\
Iteraciones y transitorio & \rule{0pt}{3.5ex} \\
Barrido: parámetro, intervalo, \(N\), observable y continuación & \rule{0pt}{5.0ex} \\
Lyapunov: burn-in, tiempo final e intervalo QR & \rule{0pt}{4.5ex} \\
Espectro: método, descarte y límites de frecuencia & \rule{0pt}{4.5ex} \\
Cuenca: región, malla, plano, horizonte y clases & \rule{0pt}{5.0ex} \\
Proyección y estilo visual & \rule{0pt}{4.5ex} \\
Estado, razón, rangos y advertencias & \rule{0pt}{5.0ex} \\
Salida esperada antes de calcular & \rule{0pt}{4.5ex} \\
Observación obtenida & \rule{0pt}{6.0ex} \\
Cambio exacto de la prueba de refinamiento & \rule{0pt}{5.0ex} \\
Comparación base frente a refinada & \rule{0pt}{6.0ex} \\
Qué permite sostener la corrida & \rule{0pt}{5.0ex} \\
Evidencia adicional requerida & \rule{0pt}{5.0ex} \\
Archivos y relación entre ellos & \rule{0pt}{5.0ex} \\
SHA-256 o suma de comprobación, si se calculó & \rule{0pt}{3.5ex} \\
\bottomrule
\end{longtable}

\begin{cajaverde}{Cierre de una receta}
Una receta está completa cuando la pregunta inicial puede responderse con una
comparación explícita entre corrida base y refinada, los archivos pueden
relacionarse sin depender de la memoria de quien los produjo y la conclusión
usa un verbo proporcional a la evidencia: ``muestra'', ``es compatible con'',
``respalda numéricamente'' o ``requiere pruebas adicionales''.
\end{cajaverde}

\begin{cajaalerta}{Vocabulario de prudencia científica}
Evite convertir automáticamente ``trayectoria acotada'' en ``atractor'',
``figura irregular'' en ``caos'', ``coexistencia'' en una clasificación global o
``resultado reproducido una vez'' en ``resultado robusto''. Cada una de esas
transiciones necesita hipótesis, diagnósticos y controles adicionales.
\end{cajaalerta}
